\chapter{Scripts and Markdown in Panels}\label{ch:panels}

A Panel is Grasshopper's text box. People keep two kinds of text in Panels: scripts, fed to a script component's
\emph{script} input, and notes. Grasshopper shows both as centered plain text. miniVivisection shows scripts as
code, edits them in place, and renders notes as Markdown.

\section{Scripts in Panels shown as code}\label{sec:panelcode}

A Panel holding a Python or C\# script shows it in Grasshopper's monospace font, with line numbers and syntax
colors, unwrapped. A long line runs on past the Panel's edge; a scroll bar along the bottom (drag it, or click the
track) brings it into view, next to the Panel's own scroll bar on the right.

\shot{07_script_panel_as_code}{Lesson GH~9: the \emph{Twist script} Panel holds the script the \emph{Twist} component
  runs. It is shown as code, with line numbers and syntax colors.}{fig:panelcode}

\paragraph{When a Panel counts as a script.} When its first line is RhinoCode's marker (\code{\#! python 3},
\code{\#! python 2}, \code{// \#! csharp}); when it is connected to a script component's \emph{script} input (that
component's language); or when its lines are plainly Python or C\#: imports, \code{def} and \code{class} lines,
\code{using} lines, statements ending in \code{;}. Notes stay text. Setting: \setting{Scripts in Panels shown as code}.

\section{Editing a script Panel in place}\label{sec:paneledit}

\textbf{Double-click} a script Panel to edit it right there, looking the same:

\begin{xltabular}{\linewidth}{@{}P{0.32\linewidth}L@{}}
\toprule
mouse, \key{Shift}+arrows       & caret and selection \\
\keys{Ctrl+Z}, \keys{Ctrl+Y}    & undo, redo \\
\keys{Ctrl+C}, \keys{Ctrl+X}, \keys{Ctrl+V} & clipboard \\
\key{Enter}                     & keeps the indentation; indents after a \code{:} in Python, a \code{\{} in C\# \\
\key{Tab}, \keys{Shift+Tab}     & indent, outdent \\
\keys{Ctrl+/}                   & comments lines in or out \\
wheel, \keys{Shift+wheel}       & scroll, scroll sideways \\
click elsewhere, \keys{Ctrl+Enter}, \keys{Ctrl+S} & keep the changes: one step of Grasshopper's Undo \\
\key{Esc}                       & drop them \\
\bottomrule
\end{xltabular}

Its right-click menu has \ui{V° Edit the script here}, and \ui{V° Edit the script in the Script Editor}: a working file
(\file{.py} or \file{.cs}) with completion and everything else, saved back into the Panel. A Panel that feeds a script
component's \emph{script} input opens as that component's script: the editor knows the component's inputs as variables,
and its \ui{Run} and \ui{Debug} compute the component in Grasshopper -- with its inputs, the editor's text taken into the
Panel first --, not as a Rhino script, which has none. \keys{Shift+double-click}
opens Grasshopper's own text box, without wrapping, in the monospace font, \key{Tab} as four spaces;
\keys{Ctrl+double-click} its dialog.

\section{After Ctrl+Alt+B stopped a Panel's script}\label{sec:panelstop}

When \keys{Ctrl+Alt+B} stopped a component that runs a Panel's script, the Panel opens for editing at the line where it
stopped: the line marked, and an explanation over it, top right: what happened, the run's data, the values then,
the call stack. Click it to put it away. The Panel keeps the line marked, with a line about it along its bottom,
until its text changes. When the Panel isn't on screen, the Script Editor shows it all.

\shot{08_panel_edited_at_the_stop}{The script Panel open for editing at the stop: line 8 marked, and the explanation:
  \emph{Stopped in 'Twist' at line 8, run 1 (Ctrl+Alt+B)}, the values \code{n = 12, angle = 5.0, i = 7}, and what to do.}{fig:panelstop}

\section{Scripts from a component's script input}\label{sec:scriptinput}

A script component's \emph{script} input can be fed by a Panel's text or by a file's path. Grasshopper then gives the
component no editor of its own. miniVivisection opens the script in the Script Editor instead, as a file:

\begin{description}
  \item[A file's path] opens that file. Saving it, or closing the Script Editor, recomputes the component.
  \item[A Panel's text] opens as a working file. Saving, or closing the Script Editor, writes the text back into the Panel,
        as an undo step, and the component recomputes. If the Panel changed meanwhile, the save is refused: open it again.
  \item[Text from anywhere else] opens as a copy to read.
\end{description}

Open it with \menu{right-click the component\>V° Edit the script input in the Script Editor}, or by double-clicking the
component. It always shows the Panel's text as it is now; when the Panel changed while the editor had changes that weren't
saved, those are kept in a file next to the working copy (\vn{Log} names it). After \keys{Ctrl+Alt+B} stopped such a
component, its script opens by itself at the line where it stopped.

\paragraph{Debugging it.} \ui{V° Debug the script input} in the component's menu runs it once in the debugger, pausing at
its breakpoints, or at its first line when it has none. Under \emph{Always debug} every run does, as long as the
file holds what Grasshopper runs (a Panel's working copy equal to its text).

\section{A component with a script input}\label{sec:libraryinput}

Grasshopper switches its own editor off for such a component. \textbf{Double-click it}, and its \emph{library} input's
files (Python or C\#, up to 20) and its script open in the Script Editor, the script's tab in front. Its menu's
\ui{Open Script Editor} does the same, and \ui{Export Script} saves its script to a file (Grasshopper grays both out
for such a component). Any script component with a library input has \ui{V° Open the library's files in the Script
Editor} in its menu.

\section{Show as}\label{sec:showas}

Every Panel's menu has \ui{V° Show as}, saved with the definition:

\begin{xltabular}{\linewidth}{@{}P{0.28\linewidth}L@{}}
\toprule
\ui{Autodetect (code or Markdown)} (the default) & code when it is a Python or C\# script; rendered when it is plainly Markdown, that is two kinds of strong signs such as \textbf{bold}, \code{code}, links, lists, quotes or code fences, or one of them and a \code{\#} heading; else as Grasshopper shows it \\
\ui{Python code}, \ui{C\# code} & any Panel of its own text shown and edited as code, also one nothing marks as a script \\
\ui{Markdown (rendered)} & rendered (below) \\
\ui{Plain text} & never code or Markdown: shown from the top left, as code is, in Grasshopper's font and wrapped as its \emph{Wrap Items} says (Grasshopper centers a Panel's text), with the same scroll bars \\
\bottomrule
\end{xltabular}

\section{Markdown in Panels}\label{sec:panelmarkdown}

A Markdown Panel shows headings, \textbf{bold}, \emph{italic}, \code{code}, links, lists, quotes, code blocks, tables
(as fixed-width text) and rules, laid out in the Panel's width.

\shot{09_markdown_panel_links}{Lesson GH~11: the lesson Panel and the \emph{Notes} Panel rendered as Markdown. The mouse
  is over a link; its tooltip says where it goes and how to follow it.}{fig:panelmarkdown}

\begin{itemize}
  \item \textbf{Double-click a word}, and the Markdown source opens in the Panel's editor with the caret at that very place.
        \key{Enter} goes on with a list or a quote. Click elsewhere and it is rendered again.
  \item \ui{V° Edit the Markdown here} in its menu opens the source in the Panel too.
        \ui{V° Edit the Markdown in the Script Editor} opens it in the Script Editor, which shows Markdown too; saving
        writes it back into the Panel.
  \item \textbf{Click a link} to follow it -- or \keys{Ctrl+click}, as \setting{Markdown links: follow with} says
        (\cref{ch:links}). The pointer becomes a hand, and a tooltip says where it goes and how; a double-click still edits.
  \item \textbf{Click a task box} (\code{- [ ]}) the same way: the task is checked or unchecked in the Panel's text, an
        undo step. A Panel fed by another isn't changed: its text comes from upstream.
\end{itemize}

\paragraph{The look.} \setting{Markdown style for Grasshopper}: a Markdown style as for Rhino text (\see{sec:mdstyle}),
its sizes in em of the Panel's font size -- \file{styles\textbackslash markdown_grasshopper.css} by default, with a line
for the quote's bar and shaded code. A change of the file shows at once. V°Text draws Rhino text, so it takes the
Rhino file's style instead.

\paragraph{The parser.} \setting{Markdown parser}: \emph{Rhino} (the default) is Markdig, Rhino's own, as in the Script
Editor's preview: CommonMark's rules, nested lists and quotes, task lists, footnotes, tables only with
their header rule. \emph{miniVivisection} is the plug-in's simpler parser. The same parser serves \code{MarkdownText} and
V°Text. Setting \setting{Markdown in Panels} off: Markdown Panels show their text as Grasshopper does.

\section{Panels fed by another}\label{sec:fedpanels}

A Panel whose input is connected shows its data, one item a line, the same way, when it opens like Markdown (a
heading, a list) or starts with a script's marker line; \ui{V° Show as} in its menu decides otherwise. It isn't edited
in place: its data comes from upstream.

\begin{lesson}{lessons GH~9 to GH~12}
\lessonfile{gh9_script_from_a_panel}: a script from a Panel, saved back. \lessonfile{gh10_script_from_a_file}: a script
file, \keys{Ctrl+Alt+B} opens it at its line. \lessonfile{gh11_markdown_notes}: Markdown notes, editing where you
double-click, links. \lessonfile{gh12_markdown_parser}: Markdig's task boxes, nested lists, footnotes.
\end{lesson}
